\begin{frame}{Core mechanism: assessment as development}

\begin{exampleblock}{\faListCheck~One rubric, many iterations}
Same assessment rubric reused across all modules and the capstone

\vspace{.3em}

\small\textit{Standards fixed; performance expectations scaled during learning}
\end{exampleblock}

\begin{columns}[T,onlytextwidth]
\begin{column}{.48\textwidth}
\begin{block}{\faClipboardCheck~Assessment criteria}
\begin{itemize}
    \item analytical soundness
    \item transparency
    \item reproducibility
    \item stakeholder relevance
\end{itemize}
\end{block}
\end{column}

\begin{column}{.48\textwidth}
\begin{alertblock}{\faForward~Developmental shift}
\centering
Feedback is \emph{carried forward} \\
---  \textit{not} reset between tasks
\end{alertblock}
\end{column}
\end{columns}

\end{frame}
\note{
A key mechanism is how expectations are made explicit.  

From the beginning, students are shown the \emph{final rubric} --  
so they know what high-quality work looks like.  

However, during the learning phase,  
assessment is \emph{scaled}.  

So instead of grading against the final standard,  
work is normalized to what counts as \emph{``very good''}.  

This signals clearly:  
this is what is expected for the final project.  

It also helps students prioritize their effort --  
they can identify what needs to be improved and revisit earlier work.  

So the rubric is not just used for grading --  
it becomes a tool for \emph{guiding development over time}.

This is important because it shifts effort  
from producing isolated outputs  
to improving work toward something that could actually be used in practice.
}
